\chapter{Project Practice and Durable Task Tracking}
\index{task tracking}\index{source audit}\index{revision pin}
\label{ch:project-practice}

CNA's repository is rich in plans, handoffs, audits, remediation ledgers, and integration records.
That abundance is valuable and dangerous: an old plan can explain why code exists while being
wrong about whether it exists today.  Durable practice is therefore not ``trust the tracking
document.''  It is preserving a chain from intent to source, test, artifact, and current status.

\section{Documents have roles, not equal authority}

The root \texttt{NEXT.md} is a session and integration handoff. At alpha.1 it opens with the
2026-08-18 merge and compiled-effect follow-up, while thousands of lines of earlier renderer,
platform, C API, glTF, and remediation history remain below it. It is operational history, not a
normalized single source of current feature truth.

\texttt{CHECKLIST.md} defines repeatable port rules: API surface comes from XNA reference
assemblies, behavior from FNA, extensions require CNAEXT, and each public member needs tests or an
explicit host-dependent exception.  It is a policy document; compliance still has to be checked
against code and gates.

\texttt{AUDIT.md} is a large historical class table with a current C-ABI tracking banner above
older per-class rows. It contains valuable provenance and many dated claims; current banner text
does not automatically refresh older rows. Treat it as a claim corpus that must be checked against
the tag, not a normalized current API inventory. The same warning applies to source-audit files
generated before modularization.

\section{Subsystem plans preserve causality}

The CNA root contains 62 \texttt{plan\_*.md} files at the alpha.1 tag, not a small fixed set. Their
numbered tasks record rationale, experiments, owner decisions, defects discovered while doing the
work, and follow-up boundaries.  A stable identifier such as \texttt{GLTF-245} or
\texttt{REMED-GFX-188} is useful because commits, tests, and later tasks can point back to the same
decision.

A done marker is still a claim.  Several audits found tasks marked complete because code existed,
not because their acceptance case ran; others found plan banners saying nothing was implemented
above sections describing completed phases.  Before quoting a task, check its named path, test,
commit, and present behavior.  If any one moved, update the forward pointer rather than erasing the
old identifier.

When a task expands across independent renderers or failure modes, split it.  The parent becomes a
tracking pointer and each child receives its own acceptance evidence.  This keeps a partial fix
from turning a broad row green and lets controls finish independently.

\section{Integration records need immutable inputs}

The renderer-integration campaign demonstrates a stronger form of handoff.  It records every
feature lane's exact SHA, derives merge order from path interaction, preserves signed no-fast-
forward merges, recomputes registry arithmetic, and names conflicts whose correct result could not
be taken from either lane alone.  A later reader can reconstruct which trees were combined and why
the result differs from each parent.

For work of this scale, a useful checkpoint packet contains:

\begin{itemize}
  \item base and input SHAs, branch roles, and proof no hidden worktree state was consumed;
  \item path-overlap and ordering rationale;
  \item exact configure/build/test commands with renderer and host identity;
  \item machine-readable results, skip counts, logs, screenshots, and generated manifests;
  \item newly found defects separated from the requested change;
  \item a clean-tree/diff check and a handoff naming the next bounded step.
\end{itemize}

A prose statement that ``all lanes passed'' cannot replace this packet, especially when one lane
was only compile-checked or one renderer ran under a substitute runtime.

\section{Executable ledgers resist prose drift}

The glTF conformance work keeps known defects in both documentation and tests, then a meta-test
asserts that the two sets agree.  Skia validators compare parity ledgers and decisions with live
interfaces and enums.  Renderer identity gates compare the CMake and C\texttt{++} registries.
Module probes inspect link closures.

These mechanisms encode a general rule: if a plan row names a finite set that code can derive,
derive and compare it.  Check row uniqueness, accepted status vocabulary, set equality, path
existence, and non-empty parsing.  Keep negative fixtures and mutation so the validator can prove it
goes red.

Executable documents can rot too.  Four wired validators at the pin read paths removed by
modularization.  A gate that cannot find its input is not current evidence.  Self-validation must
include its input mapping, and a physical source move should update implementation, tests,
validators, and cited docs atomically.

\section{SemVer and phase tracking coexist}

The first tagged release adds a second durable coordinate without replacing task history.
Product version authority begins in the root CMake numeric version plus its prerelease field;
configuration generates \path{CNA/Version.hpp}, and release preparation synchronizes Doxygen,
the changelog, and the \texttt{v0.1.0-alpha.1} tag. \texttt{CNA::getVersionString()} and the
\texttt{CNA\_VERSION\_*} macros expose that product version to C++ consumers.

The experimental native C package has a separate ABI version in \path{CNA/C/abi.h}. Its package
compatibility and release gate must move according to C ABI rules, not because the C++ product tag
changed. A release record is therefore complete only when it names both coordinates where the C
surface is in scope. Alpha.1 remains pre-1.0: SemVer identifies the snapshot and communicates
prerelease status; it does not promise a stable 1.0 API.

The release checklist in \path{docs/releasing.md} treats source, changelog, generated version API,
tag, and verification evidence as a coordinated operation. Phase/task identifiers remain the
causal engineering ledger underneath that snapshot. A task answers \emph{why and how}; the tag
answers \emph{which immutable product state}.

\section{sharp-runtime's database makes scope measurable}

The sharp-runtime sibling stores coarse classifications and concrete tickets in
\texttt{plan.sqlite3}.  At pinned commit \texttt{f827a6c5}, the \texttt{task} table contains exactly
16,201 rows: 1,082 \texttt{ported}, 14,979 \texttt{ignored}, and 140 \texttt{ignore}; 118 rows also
carry \texttt{outofscope=1}.  The \texttt{ticket} table contains 2,183 work items with descriptions,
acceptance criteria, priorities, estimates, status, and optional validation commands.

The denominator matters more than the apparent percentage.  The database deliberately classifies
vast areas of .NET that sharp-runtime will never implement.  ``1,082 ported'' means rows with a
working C\texttt{++} counterpart under this schema; it is not a claim of 6.7\% behavioral parity
with the whole BCL.  Doxygen \texttt{@note Status:} tags are secondary human hints and can drift
from the database.

The query should travel with the number:

\begin{verbatim}
SELECT status, COUNT(*) FROM task GROUP BY status;
SELECT COUNT(*) FROM ticket;
\end{verbatim}

This makes a future change visible and prevents an old book number from becoming an unexplained
project constant.

\section{Separate observations, inferences, and decisions}

A robust task record distinguishes three things.  An observation is directly reproducible:
``this shader compiler path imports a missing module.''  An inference connects evidence:
``modularization left the workflow validator stale.''  A decision sets policy:
``the release gate must fail closed on a missing manifest.''

Mixing them causes plans to outlive their premises.  A decision can remain valid after the cited
bug moves; an observation can become false after a fix; an inference may need revisiting when a
new control run appears.  Marking the category makes later re-audit cheaper.

\section{Handoff is part of implementation}

An effective handoff states what is complete, what was verified, what remains, which repositories
are read-only, and which actions require external approval.  It includes failed approaches whose
repetition would waste time and identifies known-stale documents that must not be used as status
sources.  It does not paste an ever-growing transcript into one file.

The final authority order is practical: current source and generated configuration; executable
tests and retained artifacts; accepted integration/remediation records; then plans, audits,
READMEs, and comments as context.  Comments often preserve the best explanation and the worst
status.  Durable project practice keeps both facts visible by linking explanation to something a
future maintainer can rerun.
